\subsection{(Alternative definition of $\sigma$-blocking)}

Revisit Definition~\ref{def:sigma_blocking} of $\sigma$-blocking. In the literature, one also encounters
a slightly different notion with the same name. In this optional section, we will clarify the relationship
between the two. For the rest of the course, we will keep using Definition~\ref{def:sigma_blocking}.
\begin{Def}\label{def:sigma_blocking_2}
  Let $G=\langle J,V,E,L \rangle$ be a CDMG, $C \ins J \cup V$ a subset of nodes and $\pi$ a walk (or path) in $G$:
  \[ \pi =\lp  v_0 \sts \cdots \sts v_n \rp.  \]
\begin{enumerate}
  \item
    We say that the walk (or path) $\pi$ is \emph{$C$-$\sigma'$-blocked} or \emph{$\sigma'$-blocked by $C$} if either:
      \begin{enumerate}[label=\roman*.)]
          \item $v_0 \in C$ or $v_n \in C$ or:
          \item there are two adjacent edges in $\pi$ of one of the following forms:
              \[\begin{array}{rccl}
  \text{ left chain: } & v_{k-1} \ot v_k \ots v_{k+1} & \text{ with } & v_k \in C \land v_k \notin \Sc^G(v_{k-1}),\\
  \text{ right chain: } & v_{k-1} \sto v_k \to v_{k+1} & \text{ with } & v_k \in C \land v_k \notin \Sc^G(v_{k+1}),\\
  \text{ fork: } & v_{k-1} \ot v_k \to v_{k+1} & \text{ with } & v_k \in C \land v_k \notin \Sc^G(v_{k-1}) \cap \Sc^G(v_{k+1}),\\
                \text{ collider: } & v_{k-1} \sto v_k \ots v_{k+1} & \text{ with } & v_k \notin \Anc^G(C).
              \end{array}\]
      \end{enumerate}
  \item We say that the walk (or path) $\pi$ is \emph{$C$-$\sigma'$-open} if it is not $C$-$\sigma'$-blocked. 
\end{enumerate}
\end{Def}
The only difference is in the condition for a collider node $v_k$: to $\sigma$-block
a walk one needs $v_k \not\in C$, whereas to $\sigma'$-block a walk one needs $v_k \not\in \Anc^G(C)$.

The following lemma will be convenient to relate the notions.
\begin{Lem}\label{lem:replace_walk}
Let $G = \langle V, J, E, L \rangle$ be a CDMG, $C \subseteq V \cup J$ and $\pi = (v_0,\dots,v_n)$ be a $C$-$\sigma'$-open walk in $G$. 
Suppose $v_i \in \Sc^G(v_j)$ with $i < j$.
We can then replace the subwalk $v_i, \dots, v_j$ of $\pi$ by 
  \begin{enumerate}[label=(\roman*)]
    \item a shortest directed path $v_i \to \dots \to v_j$ in $G$ if $j = n$ or if $v_j \to v_{j+1}$ on $\pi$, or
    \item a shortest directed path $v_i \ot \dots \ot v_j$ in $G$ otherwise,
  \end{enumerate}
that is entirely within $\Sc^G(v_i)$ and such that the modified walk $\pi'$ is still $C$-$\sigma'$-open.
\end{Lem}
\begin{proof}
$\pi'$ cannot become $C$-$\sigma'$-blocked by one of the initial nodes $v_0 \dots v_{i-1}$ or one of the final nodes $v_{j+1} \dots v_n$ on $\pi'$, since these nodes occur in the same local configuration on $\pi$ and do not $C$-$\sigma'$-block $\pi$ by assumption.
Furthermore, $\pi'$ cannot become $C$-$\sigma'$-blocked through one of the nodes strictly between $v_i$ and $v_j$ on $\pi'$ (if there are any), since these nodes are all non-endpoint non-colliders that only point to nodes in the same strongly connected component on $\pi'$.
Because $\pi$ is $C$-$\sigma'$-open, $v_j \notin C$ if $j = n$ or if $v_j \to i_{j+1}$ on $\pi$. 
This holds in particular in case (i). 
Similarly, $v_i \notin C$ if $i = 0$ or $v_{i-1} \ot v_i$ on $\pi$.

In case (i), $\pi'$ is not $C$-$\sigma'$-blocked by $v_j$ because $v_j$ is a non-collider on $\pi'$ but $v_j \notin C$. 
Also $v_i$ does not $C$-$\sigma'$-block $\pi'$. Indeed, assume $v_i \ne v_j$ (otherwise there is nothing to prove).
If $i = 0$, or if $i > 0$ and $v_{i-1} \ot v_i$ on $\pi'$, then the same holds for $\pi$ and hence $v_i \notin C$; $v_i$ is then a non-collider on $\pi'$, but $v_i \notin C$. 
If $i > 0$ and $v_{i-1} \sto v_i$ on $\pi'$ then $v_i$ is a non-endpoint non-collider on $\pi'$ that does not point to a node in another strongly connected component.

Now consider case (ii). 
If $i = 0$ or $v_{i-1} \ot v_i$ on $\pi'$ then this case is analogous to case (i). 
So assume $i > 0$ and $v_{i-1} \sto v_i$ on $\pi'$. 
%If $v_i$ is an endpoint of $\pi'$, then $v_i = v_j$ and $k = n$ and therefore $v_j \notin C$, and hence $v_i$ and $v_j$ do not $C$-$\sigma'$-block $\pi'$. 
$v_i$ must be a collider on $\pi'$ (whether $v_i = v_j$ or not). 
Then on the subwalk of $\pi$ between $v_i$ and $v_j$ there must be a directed walk from $v_i$ to a collider that is in $\Anc^G(C)$, which implies that $v_i \in \Anc^G(C)$. Hence, $v_i$ will not $C$-$\sigma'$-block $\pi'$. 
Also $v_j$ cannot $C$-$\sigma'$-block $\pi'$.
Indeed, assume $v_i \ne v_j$ (otherwise there is nothing to prove). 
Since $v_j \ots v_{j+1}$ on $\pi'$, $v_j$ is a non-endpoint non-collider on $\pi'$ that does not point to a node in another strongly connected component.
\end{proof}
With the help of this lemma, we can prove that Definition~\ref{def:sigma_blocking_2} indeed serves as an alternative to 
Definition \ref{def:sigma_blocking}.
\begin{Prp}\label{prp:sigma_opens}
  Let $G = \langle J, V, E, L \rangle$ be a CDMG. For $C \subseteq J \cup V$, and $i, j \in J \cup V$, the following are
  equivalent:
  \begin{enumerate}
    \item there exists a $C$-$\sigma$-open walk between $i$ and $j$ in $G$;
    \item there exists a $C$-$\sigma'$-open walk between $i$ and $j$ in $G$;
    \item there exists a $C$-$\sigma'$-open path between $i$ and $j$ in $G$.
  \end{enumerate}
\end{Prp}
\begin{proof}
$1 \implies 2$: Suppose there exists a $C$-$\sigma$-open walk between $i$ and $j$.
%Then all non-colliders are either not in $C$, or only point towards nodes in the same strongly connected component;
%all colliders are in $C$. 
All colliders are in $C$, and hence in $\Anc^G(C)$. So this same walk is also $C$-$\sigma'$-open. 

$2 \implies 1$: Suppose there exists a $C$-$\sigma'$-open walk between $i$ and $j$. 
%Then all non-colliders are either not in $C$, or only point towards nodes in the same strongly connected component; 
Then all colliders are in $\Anc^G(C)$. If a collider is not in $C$, we can extend the walk by replacing this collider 
$\sto v \ots$ by the walk 
$\sto v \to \dots \to c \ot \dots \ot v \ots$ consisting of a directed walk starting from $v$ towards the first
node $c$ in $C$ it encounters, and from there tracing back the same walk in opposite direction to $v$. All nodes on this
walk except for $c$ ($v$ included) are non-colliders not in $C$. Extending this walk in such a way for each
collider not in $C$, we obtain an extended walk that is $C$-$\sigma$-open.

$2 \implies 3$:
%It suffices to show that if there exists a $C$-$\sigma$-open walk between $i$ and $j$ in $G$, there exists a $C$-$\sigma'$-open path between $i$ and $j$ in $G$. 
Let $\pi = (v_0,\dots,v_n)$ be a $C$-$\sigma'$-open walk between nodes $v_0$ and $v_0$ in $G$. 
If a node $w$ occurs more than once on $\pi$, let $v_i$ be the first node on $\pi$ and $v_j$ be the last node on $\pi$ that are in $\Sc^G(w)$.
We now use Lemma~\ref{lem:replace_walk} to construct a new walk $\pi'$ from $\pi$ by replacing the subwalk between $v_i$ and $v_j$ of $\pi$ by a particular directed path in $\Sc^G(w)$ between $v_i$ and $v_j$ in such a way that $\pi'$ is still $C$-$\sigma'$-open.
In $\pi'$, the number of nodes that occurs more than once is at least one less than in $\pi$, and all nodes within $\Sc^G(w)$ occur within a single segment.
This replacement procedure can be repeated until no nodes occur more than once. We have then obtained a
$C$-$\sigma'$-open path between $v_0$ and $v_n$.

$3 \implies 2$: trivial.
\end{proof}
Hence, it does not matter which of the two notions we use in the definition of $\sigma$-separation.
When checking $\sigma$-separation in a graph ``by hand'', we usually use the third formulation, since there are only a finite
number of paths to check. In proofs, though, it often is easier to make use of walks, since these can be concatenated into walks
(while one cannot in general concatenate two paths and again obtain a path).


\subsection{Acyclifications}

It is possible to reformulate the notion of $\sigma$-separation in terms of $d$-separation on a modified and acyclic graph
by making use of the following construction, which will be the main tool to extend the acyclic theory for
L-CBNs to cyclic SCMs. The construction was first proposed in the context of CBNs by \cite{Spirtes94,Spirtes95}.
\begin{Def}\label{def:cdg_acyclification}
  Given a CDMG $G = \langle J, V, E, L \rangle$, we call a CADMG $G' = \langle J, V, E', L' \rangle$
  an \emph{acyclification of $G$} if
  \begin{enumerate}[label=(\roman*)]
    \item $G'$ is acyclic;
    \item $G'$ has the same input nodes $J$ and output nodes $V$ as $G$;
    \item for any pair of nodes $\{i,j\}$ such that $i \not\in\Sc^G(j)$:
      \begin{enumerate}
      \item $i \to j \in E'$ iff there exists a node $j' \in \Sc^G(j)$ such that $i \to j' \in E$;
      \item $i \oto j \in L'$ iff there exist nodes $i' \in \Sc^G(i), j' \in \Sc^G(j)$ such that $i' \oto j' \in L$;
      \end{enumerate}\label{def:acyclification_inter_scc}
    \item for any pair of distinct nodes $\{i,j\}$ such that $i \in\Sc^G(j)$: $i \to j \in \C{E}'$ or $i \ot j \in \C{E}'$ or $i \oto j \in \C{F}'$.\label{def:acyclification_intra_scc}
  \end{enumerate}
\end{Def}

\begin{comment}
Before we state and prove its main property, we start with a little lemma.
\begin{Lem}\label{lem:acyclification_ancestors}
Let $G$ be a CDMG and $G'$ an acyclification of $G$. 
For any node $j \in J \cup V$: if $i \in \Anc^G(j)$ then there
exists $i' \in \Sc^G(i)$ such that $i' \in \Anc^{G'}(j)$.
Vice versa: if $i \in \Anc^{G'}(j)$ then $i \in \Anc^G(j)$.
\end{Lem}
\begin{proof}
If there exists a directed walk from $i$ to $j$ in $G$,
this walk can be partitioned into maximal segments that are contained within a single strongly connected component of $G$ each.
By the definition of $G'$, there exists a directed walk in $G'$ that contains the right-most node only for each segment
(keeping the same ordering of the segments). 
Its end node will be $j$, but its start node may be a node $i' \in \Sc^G(i)$ that differs from $i$.

If there exists a directed walk from $i$ to $j$ in $G'$,
then we can replace each edge $u \to v$ on this walk (if it is not present in $G$) by a directed walk
$u \to v' \to \dots \to v$ where $v' \to \dots \to v$ is entirely within $\Sc^G(v)$, by definition of $G'$. 
In this way we obtain a directed walk in $G$ from $i$ to $j$.
\end{proof}
\end{comment}
The important property of acyclifications is that they can be used to express $\sigma$-separation in a (possibly cyclic) graph
in terms of $d$-separation in an acyclification.
\begin{Prp}\label{prp:sigma_sep_as_d_sep}
  Let $G = \langle J, V, E, L \rangle$ be a CDMG and $G' = \langle J, V, E', L' \rangle$ an acyclification of $G$. 
  Then for $A, B, C \subseteq V \cup J$ (not necessarily disjoint) subsets of nodes,
  $$A \Perp^\sigma_G B \given C \iff A \Perp^\sigma_{G'} B \given C \iff A \Perp^d_{G'} B \given C$$
\end{Prp}
\begin{comment}\begin{proof}
We will show that there is a $C$-$\sigma'$-open walk between $A$ and $B \cup J$ in $G$ if and only
if there is a $C$-$\sigma'$-open walk between $A$ and $B \cup J$ in $G'$. This is equivalent to the
existence of a $C$-$\sigma$-open walk between $A$ and $B \cup J$ in $G'$ by Proposition~\ref{prp:sigma_opens}. 
Since $G'$ is acyclic, this
is in turn equivalent to the existence of a $C$-$d$-open walk between $A$ and $B \cup J$ in $G'$.

$\implies$: Suppose there is a $C$-$\sigma'$-open walk $\pi = (v_0, \dots, v_n)$ between $A$ and $B \cup J$ in $G$. 
All its colliders are in $\Anc^G(C)$
and all its non-colliders are either not in $C$, or otherwise, point only to nodes in the same strongly connected component.
Note that each edge between two nodes in different strongly connected components in $G$ is also present in $G'$.
Edges between two nodes in the same strongly connected component, however, may not be present in $G'$.
Therefore, we will replace these edges with walks in $G'$. 
We may also need to replace colliders in $\Anc^G(C)$ with colliders in $\Anc^{G'}(C)$.
Consider a subwalk $(v_i,\dots,v_j)$ of maximum length that is entirely contained within a strongly connected component in $G$
(with possibly $i=j$).
We distinguish different cases and show for each case
how this subwalk can be replaced by a subwalk in $G'$.
\begin{enumerate}[label=(\roman*)]
  \item 
$\sto v_i \cdots v_j \ots$: the subwalk between $v_i$ and $v_j$ has to contain a collider, say $w$,
which must be in $\Anc^G(C)$ since the walk between $i$ and $j$ is $C$-$\sigma'$-open.
By Lemma~\ref{lem:acyclification_ancestors}, there must exist a node $w' \in \Sc^G(w)$
that is in $\Anc^{G'}(C)$. Thus we can replace this subwalk by
$\sto w' \ots$ in $G'$ such that $w'$ becomes a collider in $\Anc^{G'}(C)$.

\item
$(\ot) v_i \cdots v_j \ots$:\footnote{We put parentheses around the first directed edge to indicate that
this case also applies if $v_i$ is an endnode, i.e., if $i = 0$.} 
here $v_i$ is a non-collider pointing to another strongly-connected component or $v_i$ is an endnode,
and in both cases, $v_i \notin C$. Therefore, we can replace the subwalk by $(\ot) v_i \ots$ in $G'$,
such that $v_i$ becomes a non-collider not in $C$.

\item
$\sto v_i \cdots v_j (\to)$: analogous to the previous case, we can replace it by $\sto v_j (\to)$ in $G'$,
such that $v_j$ becomes a non-collider not in $C$.

\item
$(\ot) v_i \cdots v_j (\to)$: 
$v_i,v_j$ are both not in $C$ by assumption. If $i = j$, we replace this subwalk by $(\ot) v_i (\to)$ 
such that $v_i$ becomes a non-collider not in $C$. If $i < j$, we replace this subwalk by
$(\ot) v_i \sts v_j (\to)$ with $v_i \sts v_j$ any edge connecting $v_i$ and $v_j$ in $G'$,
such that both $v_i$ and $v_j$ become non-colliders not in $C$.
\end{enumerate}
By replacing all maximal subwalks of the original walk $\pi$ that are contained within a single strongly connected
component of $G$ in this way, we obtain a walk in the acyclification $G'$ that is $C$-$\sigma'$-open by construction.
Note that the modified walk has the same endpoints ($v_0$ and $v_n$) as the original walk.

$\impliedby$: Suppose there is a $C$-$\sigma'$-open walk $\pi$ between $A$ and $B \cup J$ in $G'$.
All its colliders are in $\Anc^{G'}(C)$ (and hence in $\Anc^G(C)$ by Lemma~\ref{lem:acyclification_ancestors}), and
all its non-colliders are not in $C$. 
We will construct a walk $\pi'$ in $G$ with the same endpoints as $\pi$ that is $C$-$\sigma'$-open.

First, we replace all non-trivial maximum subwalks $v_i \sts \dots \sts v_j$ on $\pi$ consisting entirely
of edges between nodes in the same strongly connected component in $G$ by
a directed path $v_i \to \dots \to v_j$ or $v_i \ot \dots \ot v_j$ 
entirely in $\Sc^G(i)$ in the same way as in Lemma~\ref{lem:replace_walk}.
By the same reasoning as in the proof of the lemma, both $v_i$ and $v_j$ will become
either colliders in $\Anc^G(C)$, or non-colliders not in $C$, or non-colliders in $C$
that only point to a node in the same strongly connected component of $G$.

Next, for any directed edge $i \to j$ on $\pi$ with $j \notin \Sc^G(i)$, there
must be a $j' \in \Sc^G(j)$ such that $i \to j'$ is present in $G$, and hence there must be a directed path
$j' \to \dots \to j$ entirely in $\Sc^G(j)$ such that we can use
$i \to j' \to \dots \to j$ as replacement in $G$ of the edge $i \to j$ on $\pi$.
Similarly, for any bidirected edge $i \oto j$ on $\pi$ with $j \notin \Sc^G(i)$, there
must be $i' \in \Sc^G(i)$ and $j' \in \Sc^G(j)$ such that $i' \oto j'$ is present in $G$, 
and hence there must be a walk $i \ot \dots \ot i' \oto j' \to \dots \to j$ in $G$, where
$i \ot \dots \ot i'$ is entirely in $\Sc^G(i)$ and $j' \to \dots \to j$ is entirely
in $\Sc^G(j)$, that we can use as replacement in $G$ of the edge $i \oto j$ on $\pi$.
The new nodes introduced on $\pi'$ in these replacements are non-colliders
that only point to nodes in the same strongly connected component. 
The endpoints of the replacement paths do not change their status:
if they were colliders in $\Anc^{G'}(C) \subseteq \Anc^{G}(C)$ on $\pi$ they still are 
on $\pi'$, and if they were non-colliders not in $C$ on $\pi$ they still are on $\pi'$.

Hence we have constructed a walk $\pi'$ in $G$ with the same endpoints as $\pi$ that is $C$-$\sigma'$-open.
\end{proof}\end{comment}
\begin{proof}
We will show that there is a $C$-$\sigma$-open walk between $A$ and $B \cup J$ in $G$ if and only
if there is a $C$-$\sigma$-open walk between $A$ and $B \cup J$ in $G'$. Since $G'$ is acyclic, this
is in turn equivalent to the existence of a $C$-$d$-open walk between $A$ and $B \cup J$ in $G'$.

$\implies$: Suppose there is a $C$-$\sigma$-open walk $\pi = (v_0, \dots, v_n)$ between $A$ and $B \cup J$ in $G$. 
All its colliders are in $C$
and all its non-colliders are either not in $C$, or otherwise, point only to nodes in the same strongly connected component.
Note that each edge between two nodes in different strongly connected components in $G$ is also present in $G'$.
Edges between two nodes in the same strongly connected component, however, may not be present in $G'$.
Therefore, we will replace these edges with walks in $G'$. 
Consider a subwalk $(v_i,\dots,v_j)$ of maximum length that is entirely contained within a strongly connected component in $G$
(with possibly $i=j$).
We distinguish different cases and show for each case how this subwalk can be replaced by a subwalk in $G'$.
\begin{enumerate}[label=(\roman*)]
  \item 
$\sto v_i \cdots v_j \ots$: the subwalk between $v_i$ and $v_j$ has to contain a collider, say $w$,
which must be in $C$ since the walk between $v_i$ and $v_j$ is $C$-$\sigma$-open.
We can replace this subwalk by $\sto w \ots$ in $G'$ such that $w$ becomes a collider in $C$.

\item
$(\ot) v_i \cdots v_j \ots$:\footnote{We put parentheses around the first directed edge to indicate that
this case also applies if $v_i$ is an endnode, i.e., if $i = 0$.} 
here $v_i$ is a non-collider pointing to another strongly-connected component or $v_i$ is an endnode,
and in both cases, $v_i \notin C$. Therefore, we can replace the subwalk by $(\ot) v_i \ots$ in $G'$,
such that $v_i$ becomes a non-collider not in $C$.

\item
$\sto v_i \cdots v_j (\to)$: analogous to the previous case, we can replace it by $\sto v_j (\to)$ in $G'$,
such that $v_j$ becomes a non-collider not in $C$.

\item
$(\ot) v_i \cdots v_j (\to)$: 
$v_i,v_j$ are both not in $C$ by assumption. If $i = j$, we replace this subwalk by $(\ot) v_i (\to)$ 
such that $v_i$ becomes a non-collider not in $C$. If $i < j$, we replace this subwalk by
$(\ot) v_i \sts v_j (\to)$ with $v_i \sts v_j$ any edge connecting $v_i$ and $v_j$ in $G'$,
such that both $v_i$ and $v_j$ become non-colliders not in $C$.
\end{enumerate}
By replacing all maximal subwalks of the original walk $\pi$ that are contained within a single strongly connected
component of $G$ in this way, we obtain a walk in the acyclification $G'$ that is $C$-$\sigma$-open by construction.
Note that the modified walk has the same endpoints ($v_0$ and $v_n$) as the original walk.

$\impliedby$: Suppose there is a $C$-$\sigma$-open walk $\pi'$ between $A$ and $B \cup J$ in $G'$.
All its colliders are in $C$, and all its non-colliders are not in $C$. 
We will construct a walk $\pi$ in $G$ with the same endpoints as $\pi'$ that is $C$-$\sigma$-open.

Consider a non-trivial subwalk $(v_i,\dots,v_j)$ on $\pi'$ of maximum length that is entirely contained within a strongly connected component of $G$. 
This subwalk may not be present in $G'$.
We distinguish different cases and show for each case how this subwalk can be replaced by a subwalk in $G$.
\begin{enumerate}[label=(\roman*)]
  \item 
$\sto v_i \cdots v_j \ots$: the subwalk between $v_i$ and $v_j$ has to contain a collider, say $w$,
which must be in $C$ since the walk between $v_i$ and $v_j$ is $C$-$\sigma$-open, and must be in
$\Sc^G(v_i) = \Sc^G(v_j)$ by assumption. We can replace this subwalk by
$\sto v_i \to \dots \to w \ot \dots \ot v_j \ots$ in $G$,
with possibly $v_i = w$ and possibly $w = v_j$, 
with all nodes in $\Sc^G(v_i)$. 
Note that the modified walk remains $C$-$\sigma$-open.

\item
$(\ot) v_i \cdots v_j \ots$:
here $v_i$ is a non-collider pointing to another strongly-connected component or $v_i$ is an endnode,
and in both cases, $v_i \notin C$. 
We can replace this subwalk by a shortest directed walk
$(\ot) v_i \ot \dots \ot v_j \ots$ in $G$ with all nodes in $\Sc^G(v_i)$.
Note that the modified walk remains $C$-$\sigma$-open.
%The modified walk is not $C$-$\sigma$-blocked by $v_i$ because $v_i$ is a non-collider not in $C$.
%Also $v_j$ does not $C$-$\sigma$-block the modified walk, because either
%$v_i \ne v_j$ or otherwise $v_j$ is a non-collider pointing only to nodes in the same strongly connected component.

\item
$\sto v_i \cdots v_j (\to)$: analogous to the previous case, we can replace it by $\sto v_i \to \dots \to v_j (\to)$ in $G$.

\item
$(\ot) v_i \cdots v_j (\to)$: 
$v_i,v_j$ are both not in $C$ by assumption. We can replace this subwalk by a shortest directed walk
$(\ot) v_i \to \dots \to v_j (\to)$ in $G$ with all nodes in $\Sc^G(v_i)$.
The modified walk remains $C$-$\sigma$-open.
\end{enumerate}
In each of the four cases, in the modified walk both $v_i$ and $v_j$ become
either colliders in $C$, or non-colliders not in $C$, or non-colliders in $C$
that only point to a node in the same strongly connected component of $G$.

Now, we will replace edges on $\pi'$ between two strongly connected components that are not present in $G$.
For any directed edge $i \to j$ on $\pi'$ with $j \notin \Sc^G(i)$, there
must be a $j' \in \Sc^G(j)$ such that $i \to j'$ is present in $G$, and hence there must be a directed path
$j' \to \dots \to j$ entirely in $\Sc^G(j)$ such that we can use
$i \to j' \to \dots \to j$ as replacement in $G$ of the edge $i \to j$.
Similarly, for any bidirected edge $i \oto j$ on $\pi'$ with $j \notin \Sc^G(i)$, there
must be $i' \in \Sc^G(i)$ and $j' \in \Sc^G(j)$ such that $i' \oto j'$ is present in $G$, 
and hence there must be a walk $i \ot \dots \ot i' \oto j' \to \dots \to j$ in $G$, where
$i \ot \dots \ot i'$ is entirely in $\Sc^G(i)$ and $j' \to \dots \to j$ is entirely
in $\Sc^G(j)$, that we can use as replacement in $G$ of the edge $i \oto j$.
The new nodes introduced on $\pi$ in these replacements are non-colliders
that only point to nodes in the same strongly connected component. 
The endpoints of the replacement paths do not change their status:
if they were colliders in $C$ on $\pi'$ they still are 
on $\pi$, and if they were non-colliders not in $C$ on $\pi'$ they still are on $\pi$.

Hence we have constructed a walk $\pi$ in $G$ with the same endpoints as $\pi'$ that is $C$-$\sigma$-open.
\end{proof}
